\section{HASIL DAN PEMBAHASAN}

\subsection{Akuisisi Data dan Prapemrosesan}

Proses prapemrosesan data harga historis saham menghasilkan dua parameter utama yang menjadi masukan formulasi Markowitz yaitu ekspektasi imbal hasil sederhana $\mu_i$ dan matriks kovarians $\Sigma$ yang dihitung menggunakan imbal hasil logaritmik pada jendela pelatihan 6 bulan atau $n=126$ hari perdagangan. Penggunaan imbal hasil sederhana untuk $\mu_i$ didasarkan pada sifat liniernya terhadap bobot aset, sedangkan imbal hasil logaritmik untuk $\Sigma$ dipilih karena memiliki distribusi yang lebih mendekati asumsi normalitas. Analisis melalui aproksimasi deret Taylor orde kedua menunjukkan bahwa kedua representasi menghasilkan estimasi kovariansi yang hampir identik untuk data saham harian, yakni $\text{Cov}(r_A, r_B) \approx \text{Cov}(R_A, R_B)$, karena suku koreksi berada pada orde $10^{-6}$ hingga $10^{-8}$ yang jauh lebih kecil dari suku utamanya. Pada sistem $N=2$ dengan jendela pelatihan pertama, diperoleh matriks kovarians 
\begin{equation*}
\Sigma_{N2} = \begin{pmatrix} 0{,}0915 & 0{,}0202 \\ 0{,}0202 & 0{,}0367 \end{pmatrix},
\end{equation*}
sedangkan parameter penghindaran risiko endogen $\gamma$ dihitung secara adaptif menggunakan fungsi sigmoid 
\begin{equation*}
\gamma = 1/(1+e^{\bar{\mu}_l/\bar{\sigma}_l}) \approx 0{,}2697
\end{equation*}
pada jendela pertama.

\subsection{Formulasi EPG dan Pencarian Ekuilibrium Nash}

Transformasi model Markowitz ke dalam kerangka EPG dilakukan dengan mendefinisikan fungsi potensial global sebagaimana pada persamaan (\ref{eq:global_potential_function}) yang merangkum seluruh interaksi strategis antar aset. Kemudian pencarian ekuilibrium Nash dilakukan menggunakan algoritma SBR yang bekerja melalui mekanisme pertukaran aset secara iteratif yaitu pada setiap iterasi, seluruh kandidat pertukaran antara aset yang ada di dalam dan di luar portofolio dievaluasi lalu pertukaran yang memberikan kenaikan fungsi potensial terbesar dipilih hingga tidak ada lagi pertukaran yang meningkatkan $\Phi$. Sebagai ilustrasi pada sistem $N=2$ periode Januari 2021, evaluasi seluruh kombinasi strategi mengidentifikasi bahwa keadaan $|10\rangle$ (memilih BBCA, tidak ADRO) menghasilkan utilitas sistem tertinggi sehingga ditetapkan sebagai titik ekuilibrium Nash. Konfigurasi biner yang dihasilkan EPG ini selanjutnya dikonversi menjadi \textit{bitstring} inisialisasi yang menjadi \textit{warm-start} bagi optimasi VQE.

\subsection{Pemetaan Hamiltonian dan Inisialisasi Sirkuit Kuantum}

Fungsi potensial EPG ditransformasi menjadi fungsi energi QUBO melalui relasi $E(\mathbf{x}) = -\Phi(\mathbf{x})$, yang kemudian diperluas dengan suku penalti kuadratik supaya batasan kardinalitas seleksi aset terpenuhi. Pemetaan domain spin melalui substitusi transformasi Affine dilakukan untuk mengekstraksi parameter Hamiltonian Ising berupa koefisien bias lokal $h_i$ dan parameter kopling $J_{ij}$. Sebagai contoh numerik pada sistem $N=2$ jendela pertama, Hamiltonian yang diperoleh adalah 
\begin{equation*}
    \hat{\mathcal{H}}_{N2} = 0{,}1735\hat{Z}_1 + 0{,}0931\hat{Z}_2 + 0{,}5014\hat{Z}_1\hat{Z}_2 + 0{,}2320.
\end{equation*}
Inisialisasi sirkuit kuantum kemudian dilakukan dengan merotasi keadaan dasar $|00\ldots0\rangle$ menuju \textit{bitstring} ekuilibrium Nash. Jika bitstring yang terpilih adalah $0$, maka rotasi yang digunakan adalah $0$ sedangkan jika yang terpilih adalah $1$, maka rotasi yang digunakan adalah $\pi$ rad atau $180$ derajat atau $3,14$ rad. Sehingga sirkuit memulai optimasi dari representasi kuantum keadaan yang telah terbukti memiliki utilitas finansial tinggi.

\subsection{Konvergensi Energi dan Optimasi SPSA}

Proses optimasi VQE menggunakan algoritma SPSA menunjukkan dinamika konvergensi yang berbeda secara signifikan antara skenario GT (dengan inisialisasi Nash) dan NoGT (inisialisasi acak). Pada skenario GT-VQE, konvergensi energi ekspektasi $E(\boldsymbol{\theta})$ menuju nilai minimum umumnya tercapai pada kedalaman sirkuit yang dangkal ($L=2$ hingga $L=3$), karena parameter awal yang diperoleh dari ekuilibrium Nash telah berada di wilayah dekat minimum energi sejati. Sebaliknya, pada skenario NoGT, VQE sering membutuhkan kedalaman yang lebih tinggi dan jumlah iterasi yang lebih banyak untuk mencapai konvergensi serupa, bahkan pada beberapa jendela observasi gagal menemukan keadaan dasar yang benar. Dinamika adaptasi tingkat kedalaman sirkuit per jendela simulasi ini direkam secara komprehensif, dan implikasinya terhadap stabilitas performa sistem akan dievaluasi lebih lanjut pada pembahasan evaluasi portofolio.

\begin{figure}[t!]
    \centering
    \includegraphics[width=\linewidth]{NoGT_grafik_depth_per_window_N2.png}\\
    (a)\\[0.2cm]
    \includegraphics[width=\linewidth]{GT_grafik_depth_per_window_N2.png}\\
    (b)
    \caption{Perbandingan dinamika kedalaman sirkuit untuk sistem $N=2$ antara metode VQE murni (a) dan GT-VQE (b).}
    \label{fig:depth_n2}
\end{figure}

\begin{figure}[t!]
    \centering
    \includegraphics[width=\linewidth]{NoGT_grafik_depth_per_window_N4.png}\\
    (a)\\[0.2cm]
    \includegraphics[width=\linewidth]{GT_grafik_depth_per_window_N4.png}\\
    (b)
    \caption{Perbandingan dinamika kedalaman sirkuit untuk sistem $N=4$ antara metode VQE murni (a) dan GT-VQE (b).}
    \label{fig:depth_n4}
\end{figure}

\subsection{Analisis Entropi Von Neumann}

\begin{figure*}[t!]
    \centering
    \begin{minipage}{0.48\textwidth}
        \centering
        \includegraphics[width=\linewidth]{Entropi_tinggi.png}
        \caption{Distribusi probabilitas merata yang mengindikasikan jebakan \textit{barren plateaus} (Entropi Tinggi).}
        \label{fig:entropi_tinggi}
    \end{minipage}\hfill
    \begin{minipage}{0.48\textwidth}
        \centering
        \includegraphics[width=\linewidth]{Entropi_rendah.png}
        \caption{Kolaps probabilitas pada keadaan dominan akibat inisialisasi \textit{Nash Equilibrium} (Entropi Rendah).}
        \label{fig:entropi_rendah}
    \end{minipage}
\end{figure*}

Entropi Von Neumann digunakan sebagai indikator kuantitatif untuk mendeteksi keberhasilan proses optimasi dalam menghindari fenomena \textit{barren plateaus}. Pada skenario inisialisasi acak (NoGT), nilai entropi cenderung bertahan tinggi sepanjang iterasi karena distribusi probabilitas tetap tersebar merata di seluruh basis state, sebagaimana diilustrasikan pada Gambar~\ref{fig:entropi_tinggi}. Kondisi ini secara analitik mengonfirmasi bahwa sirkuit terperangkap dalam jebakan \textit{barren plateaus}, di mana gradien ekspektasi fungsi biaya meluruh menuju nol sehingga SPSA tidak mampu menghasilkan pembaruan parameter yang efektif. Sebaliknya, implementasi strategi \textit{warm-start} berbasis ekuilibrium Nash pada skenario GT-VQE memicu keruntuhan entropi yang signifikan sejak iterasi awal, ditandai oleh pemusatan distribusi probabilitas pada satu \textit{bitstring} dominan sebagaimana ditunjukkan pada Gambar~\ref{fig:entropi_rendah}. Fenomena kolaps entropi ini memvalidasi secara fisis bahwa sistem berhasil keluar dari superposisi yang terlalu menyebar menuju keadaan kuantum yang terlokalisasi, yang secara langsung berkorelasi dengan keberhasilan menemukan konfigurasi portofolio optimal.

\subsection{Performa Finansial dan Backtesting Portofolio}

Evaluasi performa metrik simulasi pengujian historis pada periode 2021 hingga 2023 dengan mekanisme penyeimbangan kembali bulanan mengungkapkan keunggulan strategi GT-VQE secara konsisten pada kedua skala sistem. Pada sistem $N=2$, algoritma GT-VQE mencatatkan imbal hasil kumulatif sebesar $145{,}07\%$ dengan rasio Sharpe $1{,}0185$, melampaui VQE murni ($121{,}30\%$, Sharpe $0{,}9331$), strategi \textit{Equal Weight} ($97{,}88\%$, Sharpe $0{,}9764$), maupun optimasi klasik SLSQP ($51{,}34\%$, Sharpe $0{,}6736$). Keunggulan ini semakin teramplifikasi pada sistem $N=4$ yang memiliki ruang solusi lebih kompleks: VQE murni mengalami degradasi performa yang signifikan dengan imbal hasil hanya $16{,}89\%$ dan Sharpe $0{,}3592$, sedangkan GT-VQE mampu mempertahankan imbal hasil $73{,}15\%$ dengan Sharpe $1{,}0107$ dan \textit{Maximum Drawdown} terendah sebesar $16{,}68\%$. Hasil lengkap perbandingan keempat strategi disajikan pada Tabel~\ref{tab:performa}, sedangkan visualisasi kurva pertumbuhan modal untuk seluruh skenario dapat ditinjau pada Lampiran.

\begin{table}[t!]
  \centering
  \caption{Ringkasan Performa Simulasi Finansial ($N=2$ dan $N=4$)}
  \label{tab:performa}
  \resizebox{\columnwidth}{!}{%
  \begin{tabular}{lcccccc}
  \toprule
  \multirow{2}{*}{\textbf{Strategi Algoritma}} & \multicolumn{3}{c}{\textbf{Sistem $N=2$}} & \multicolumn{3}{c}{\textbf{Sistem $N=4$}} \\
  \cmidrule(lr){2-4} \cmidrule(lr){5-7}
  & \textbf{Return} & \textbf{Sharpe} & \textbf{MDD} & \textbf{Return} & \textbf{Sharpe} & \textbf{MDD} \\
  \midrule
  \textbf{GT-VQE (Hibrida)} & \textbf{145,07\%} & \textbf{1,0185} & 30,33\% & \textbf{73,15\%} & \textbf{1,0107} & \textbf{16,68\%} \\
  VQE (Tanpa GT) & 121,30\% & 0,9331 & 30,33\% & 16,89\% & 0,3592 & 20,89\% \\
  SLSQP (Klasik) & 51,34\% & 0,6736 & \textbf{15,91\%} & 12,91\% & 0,2810 & 21,12\% \\
  \textit{Equal Weight} & 97,88\% & 0,9764 & 27,49\% & 35,57\% & 0,6790 & 18,42\% \\
  \bottomrule
  \end{tabular}%
  }
\end{table}

Mekanisme inisialisasi awal berbasis ekuilibrium Nash dibuktikan memiliki pengaruh yang signifikan terhadap dinamika adaptasi kedalaman sirkuit, sebagaimana diilustrasikan pada Gambar~\ref{fig:depth_n2} dan Gambar~\ref{fig:depth_n4}. Fluktuasi kedalaman sirkuit yang cukup ekstrem teramati pada optimasi VQE murni, terutama saat diimplementasikan pada dimensi sistem permasalahan yang lebih tinggi ($N=4$), yang mengindikasikan kegagalan konvergensi pada lanskap parameter yang rumit. Sebaliknya, kestabilan operasional dipertahankan oleh metode GT-VQE pada tingkat kedalaman sirkuit yang dangkal. Hal ini menunjukan bahwa titik awal optimasi yang telah diposisikan secara strategis di sekitar area minimum energi sejati oleh kerangka EPG mampu menghasilkan peningkatan performa finansial yang signifikan dibandingkan dengan pendekatan optimasi konvensional VQE. Secara kuantitatif, perolehan performa finansial untuk GT-VQE menghasillkan perkembangan portofolio yang superior, di mana konfigurasi bobot aset yang optimal berhasil dipertahankan secara konsisten pada setiap periode \textit{rebalancing} portofolio.
